% qss-manual.tex --- A short manual for qss
%
% Copyright (C) 2026 Friedrich Dominicus
% SPDX-License-Identifier: MIT
%
% Build with LuaLaTeX (twice), or from Emacs: M-x qss-manual-build
% (doc/qss-manual-reference.el), which first writes reference.tex from
% the code.

\documentclass[
  paper=a4, fontsize=11pt, parskip=half, headings=small,
  toc=bibliography, numbers=noenddot, DIV=12
]{scrreprt}

\usepackage{fontspec}
\setmonofont{DejaVu Sans Mono}[Scale=MatchLowercase]
\usepackage[english]{babel}
\usepackage{microtype}
\usepackage[autostyle]{csquotes}
\usepackage{xcolor}
\usepackage{booktabs, longtable, array, tabularx}
\usepackage{fancyvrb}
\usepackage{listings}
\usepackage[automark]{scrlayer-scrpage}
\usepackage[hidelinks]{hyperref}

\definecolor{qssframe}{gray}{0.75}
\definecolor{qssback}{gray}{0.96}

\newcommand{\code}[1]{\texttt{#1}}
\newcommand{\file}[1]{\texttt{#1}}
\newcommand{\kbd}[1]{\texttt{#1}}
\newcommand{\qss}{\textsf{qss}}

\fvset{fontsize=\small, frame=leftline, framerule=1.5pt, rulecolor=\color{qssframe},
       xleftmargin=1em, framesep=0.8em}

\lstdefinestyle{elisp}{
  language=Lisp, basicstyle=\ttfamily\small, columns=fullflexible, keepspaces=true,
  upquote=true, showstringspaces=false, frame=leftline, framerule=1.5pt,
  rulecolor=\color{qssframe}, xleftmargin=1em, framesep=0.8em,
  commentstyle=\color{black!60}, keywordstyle={}, ndkeywordstyle={},
}
\lstset{style=elisp}

\pagestyle{scrheadings}
\KOMAoptions{headsepline=true}

\title{qss}
\subtitle{Tables, one vocabulary, presentations and contracts for Emacs\\[1ex]
  \normalsize A short manual}
\author{Friedrich Dominicus}
% Written by qss-manual-reference.el.  Do not edit.
\newcommand{\qssversion}{0.19.0}

\date{Version \qssversion}

\begin{document}

\maketitle
\tableofcontents

% ======================================================================
\chapter{Introduction}

\qss{} brings the idea of Nushell to Emacs: data as tables, and one small
vocabulary that works on all of them --- files, processes, lists of property
lists, rows of an SQLite table, and whatever Nushell itself produces. Around
this idea are a few others, older ones:

\begin{itemize}
\item on SQLite the words become SQL, as in dbplyr, and SQLite is both
  source and store;
\item what is shown stays the object it shows and can be pointed at later, as
  in CLIM: \emph{presentations} and \emph{accept};
\item the public interface is checked by \emph{contracts} with blame, as in
  Racket.
\end{itemize}

Everything is Emacs Lisp and runs inside Emacs (29.1 or later). Only what is
built into Emacs is used: \code{vtable}, \code{sqlite}, \code{cl-generic} and
advice.

This manual explains the ideas and how they fit together. The README of the
package is the short version; the docstrings (\kbd{C-h f}) are the detailed
one. The reference in chapter~\ref{chap:reference} is written from the code
itself.

\section{Installation}

\begin{lstlisting}
(package-vc-install "https://github.com/FDominicus/qss")
(require 'qss)
\end{lstlisting}

Or put the directory on \code{load-path} and \code{(require 'qss)}. Nushell is
optional; without it everything else works.

\section{A first session}

\kbd{M-x qss} opens the REPL. Type a pipe, press \kbd{RET}, and the table
appears right below:

\begin{Verbatim}
~/ qss> (files "~") (where (> :size "10mb")) (sort-by size :desc) (first 5)
5 rows  [buffer]  files ~/ | where (> :size "10mb") | sort-by size desc | first 5
name                 type  size modified         mode       user
lisp.virtualmem      file   11M 2026-06-01 14:35 -rw-rw-r-- frido
...
~/ qss> (select name size)
\end{Verbatim}

The second line starts with a word, so it continues the table above. The title
line repeats every stage of the pipe, so you always see how a table came about.
\kbd{RET} on a row opens the file, \kbd{a} offers what can be done with it, and
\code{(sqlite (qss-demo-database) "Orders")} gives a small database to try the
SQL part with.

% ======================================================================
\chapter{Tables and rows}

\section{A table is a value}

A table (\code{qss-table}) is a title, a list of columns, a list of rows and,
optionally, a way to build it again. It is an ordinary Lisp value: it can be
stored in a variable, passed on, shown.

The rows are \emph{not} turned into strings or records of strings. A row of
\code{(files)} is a \code{qss-file}, a row of \code{(processes)} a
\code{qss-process}, a row from SQLite a property list. The table only knows how
to ask a row for its fields. That is why a filtered file still opens with
\kbd{RET}, and why the same words work on all of them.

\section{The protocol of rows}

What a row is, is said by a few generic functions, with a method per type of
row:

\noindent\begin{tabularx}{\linewidth}{@{}lX@{}}
\toprule
\code{qss-columns} & the columns shown for a row \\
\code{qss-field} & the raw value of a field (filtering and sorting use it) \\
\code{qss-default-action} & what \kbd{RET} does \\
\code{qss-actions} & what \kbd{a} offers: a list of (\emph{label} . \emph{function}) \\
\code{qss-stored-fields} & fields stored by \code{into} besides the columns \\
\code{qss-presentation-type} & the type for \code{qss-accept} (chapter~\ref{chap:present}) \\
\code{qss-presentation-label} & a short text for it \\
\code{qss-file-name} & the file a row stands for, if any \\
\bottomrule
\end{tabularx}

Property lists and alists have methods already, so any list of property lists
is a table:

\begin{lstlisting}
(qss-> '((:name "Ada" :born 1815) (:name "Grace" :born 1906))
       (where (< :born 1900))
       collect)
;; => ((:name "Ada" :born 1815))
\end{lstlisting}

\section{Columns}

A column is a field name, or \code{(\emph{name} . \emph{plist})} with
\code{:type} (\code{number}, \code{size}, \code{duration}, \code{time}),
\code{:format}, \code{:align} and \code{:width}. The type decides display,
alignment and sorting: a \code{size} is bytes shown as \code{9.8k}, a
\code{duration} nanoseconds shown as \code{3hr 20min}, a \code{time} a Lisp time
shown as \code{2026-10-22 14:40}.

% ======================================================================
\chapter{The vocabulary}

\section{Sources and words}

A pipe begins with a source and goes on with words:

\begin{lstlisting}
(qss-> (sqlite (qss-demo-database) "Orders")
       (where (> :amount 100))
       (sort-by order-date :desc)
       (select order-no amount)
       collect)
\end{lstlisting}

In the REPL and in \kbd{M-x qss-eval} the same is written without \code{qss->}.
The sources are \code{files}, \code{processes}, \code{run}, \code{sqlite},
\code{tables}, \code{nu} and \code{contracts}; any list is one too. The words:

\noindent\begin{tabularx}{\linewidth}{@{}lX@{}}
\toprule
\code{(where \emph{cond})} & rows for which the condition holds (\code{filter} is the same) \\
\code{(select \emph{col}\ldots)}, \code{(reject \emph{col}\ldots)} & keep, drop columns \\
\code{(sort-by \emph{col} [:desc])} & sort; missing values last \\
\code{(first [\emph{n}])}, \code{(last [\emph{n}])} & the first, last rows \\
\code{(get \emph{col})} & the values of one column \\
\code{(group-by \emph{col}\ldots{} \emph{agg}\ldots)} & one row per group: keys, count, aggregates; \kbd{RET} shows the group \\
\code{(summarize \emph{agg}\ldots)} & one row: count and aggregates over all rows \\
\code{(join \emph{table} :on \emph{col})} & rows of both tables whose columns match \\
\code{(update-rows :\emph{col} \emph{v}\ldots)}, \code{(delete-rows)} & change, delete stored rows (section~\ref{sec:changing}) \\
\code{(each \emph{form})} & \emph{form} with \code{it} bound to each row; the results are the new rows \\
\code{(nu \emph{command})} & the table through Nushell (section~\ref{sec:nushell}) \\
\code{(to \emph{format})} & the table as text: \code{text}, \code{org}, \code{md}, \code{csv}, \code{tsv} \\
\code{(into sqlite \emph{name} \ldots)} & store the table (section~\ref{sec:into}) \\
\code{columns}, \code{count}, \code{collect}, \code{sql} & column names, number of rows, the rows as a list, the SQL \\
\bottomrule
\end{tabularx}

Every word returns a table again (or a query, on SQLite), except
\code{count}, \code{collect}, \code{to} and \code{sql}. A derived table knows
its way back to the source: \kbd{g} in its buffer reads the source again and
applies the words once more.

A column may be written with or without a colon: \code{(select name)} and
\code{(select :name)} are the same. In code the colon is the better choice,
since then one rule holds everywhere: with a colon it is a column, without it is
Lisp. In the REPL, where one types, the short form is fine.

\begin{lstlisting}
(defun big-files-of (dir size)
  (qss-display-> (files dir) (where (> :size size))
                 (select :name :size :modified) (sort-by :size :desc)))
\end{lstlisting}

\code{(qss-display-> \ldots)} is short for \code{(qss-display (qss-> \ldots))}:
the pipe, shown as a table; it returns the buffer.

\section{Groups and joins}

\begin{Verbatim}
qss> (sqlite (qss-demo-database) "Orders") (group-by status (sum amount) (max order-date))
status     count  sum-amount  max-order-date
cancelled      1       150.0  2026-10-12
paid           3       245.5  2026-10-02
               2     1299.99  2026-10-05
\end{Verbatim}

\code{group-by} takes the columns to group by and any number of aggregates:
\code{(sum \emph{col})}, \code{(avg \emph{col})}, \code{(min \emph{col})},
\code{(max \emph{col})} and \code{(count \emph{col})}, the number of values
present. They are named \code{sum-amount} and so on, or as told:
\code{(sum amount :as total)}. \code{count}, the number of rows, is always
there. As in SQL, missing values are left out of the aggregates. The groups are
sorted by their keys, missing values last, in memory and on SQLite alike; a sum
or average of sizes is a size again. \kbd{RET} on a group shows its rows.
\code{summarize} is the same without groups: one row over all rows (also over
none: then the count is 0).

\code{join} puts the rows of two tables side by side where their columns match:

\begin{Verbatim}
qss> (sqlite (qss-demo-database) "Orders")
       (join (sqlite (qss-demo-database) "Customers") :on customer-no)
       (group-by city (sum amount)) (sort-by sum-amount :desc)
\end{Verbatim}

The other table is written as a source, like the start of a pipe. \code{:on}
names the column (the same on both sides), a pair \code{(\emph{left}
\emph{right})} for different names, or a list of pairs; without \code{:on},
all columns both tables have. \code{:kind left} keeps the rows of the left
table without a partner. A column both sides have appears once if it is joined
on; otherwise the right one gets the suffix \code{-2}. Missing values never
match, as in SQL. The rows of a join are data (property lists). The options are
data like column names: in a pipe they are not evaluated.

\section{Conditions}

Columns are keywords; everything else is Lisp, variables included:

\begin{lstlisting}
(where (> :size 10000))
(where (and (>= :date since) (in :kind kinds)))
(where (like :name "%.el"))
(where (is-null :invoice-no))
\end{lstlisting}

The comparisons mean what they mean in SQL. \code{=}, \code{<} \ldots{} work on
numbers, strings and times (\code{(> :modified "2026-10-20")} works); sizes and
durations can be written as in Nushell (\code{(> :size "10kb")},
\code{(< :took "1hr")}); \code{like} ignores case; and a missing value is NULL:
comparing with it is \emph{unknown}, neither true nor false, and \code{where}
keeps only what is true. So \code{(not (= :city "Bruchsal"))} does not keep
rows without a city --- in memory exactly as in SQLite.

If the symbol \code{it} occurs, the condition is ordinary Lisp with \code{it}
bound to the row object:

\begin{lstlisting}
(where (string-prefix-p "." (qss-field it 'name)))
\end{lstlisting}

% ======================================================================
\chapter{The REPL}

\section{Input}

\begin{itemize}
\item Several forms, or one form beginning with a source, are a pipe.
\item A form beginning with a word continues the last table, like \code{\$in}
  in Nushell; so does a line beginning with \code{|}.
\item Anything else is Lisp. The last table is in the variable
  \code{qss-last}.
\item \code{!\emph{program} \emph{args}} runs a program directly, without a
  shell; its output is shown and is also the last table (one row per line, in
  the column \code{line}).
\item \code{cd \emph{dir}} changes the directory; the prompt shows it, and
  \code{files}, \code{nu}, \code{run} and \code{!} run there --- on the remote
  host in a TRAMP directory.
\end{itemize}

\begin{Verbatim}
~/ qss> cd ~/src/project
~/src/project/ qss> !git log --oneline -20
~/src/project/ qss> (where (like :line "%fix%"))
~/src/project/ qss> | (nu "length")
\end{Verbatim}

\section{Keys}

\noindent\begin{tabularx}{\linewidth}{@{}lX@{}}
\toprule
\kbd{RET} & on the input: run it; on a row: its default action (a directory: \code{cd} there) \\
\kbd{a}, right mouse button & the actions at point (chapter~\ref{chap:actions}) \\
\kbd{w} & copy the whole table as text; \kbd{C-u w} asks for the format \\
\kbd{S} & sort by the column at point \\
\kbd{e}, \kbd{d}, \kbd{u} & in a table buffer of an SQLite table: change a cell, mark a row for deletion, undo (section~\ref{sec:changing}) \\
\kbd{C-c C-v} & take a value from a table into the input, by a click \\
\kbd{C-c C-f} & take the name of a file into the input \\
\kbd{M-p}, \kbd{M-n} & history \\
\kbd{TAB} & complete words, sources, column names, file names \\
\kbd{C-c M-o} & clear the buffer \\
\code{[buffer]} & the table in a buffer of its own, where \kbd{g} reads it again \\
\bottomrule
\end{tabularx}

Long tables show the first \code{qss-repl-max-rows} rows and a
\code{[show all]} button.

% ======================================================================
\chapter{Actions at point}\label{chap:actions}

\kbd{a}, or the right mouse button, on a cell offers four groups:

\begin{enumerate}
\item the actions of the row object --- open the file, Dired, kill the
  process \ldots{} (\code{qss-actions});
\item words with the value under point: \code{(where (= :status "paid"))},
  \code{(where (/= \ldots))}, \code{(where (> :amount 120.0))} for numbers,
  \code{(where (is-null :city))} for an empty cell, the day for a time;
  then \code{sort-by}, \code{group-by}, \code{select} and \code{reject} on that
  column, and copying the value;
\item commands with typed arguments whose first argument fits the row or the
  value (section~\ref{sec:commands});
\item copying the whole table as text, Org, Markdown, CSV or TSV.
\end{enumerate}

A word continues the \emph{whole} table, not only the rows the REPL shows, and on
SQLite it continues the query, so the condition becomes SQL. In the REPL the
word is submitted as the next input --- visible and in the history --- and it
continues the table that was clicked, even an older one; what you had typed at
the prompt stays. In a table buffer the result opens in a buffer of its own.

\code{(to org)} and \kbd{w} show the cells as the table does (\code{9.8k});
\code{(to csv)} and \code{(to tsv)} carry the plain values for other programs:
sizes in bytes, durations in nanoseconds, times with seconds.

% ======================================================================
\chapter{SQLite}

\section{Queries}

On an SQLite source the pipe stays a query until the rows are needed.
\code{where}, \code{sort-by}, \code{first}, \code{count}, \code{select},
\code{reject}, \code{columns}, \code{group-by}, \code{summarize} and
\code{join} (with a query of the same file) become SQL with \code{?}
parameters; a word that cannot (a condition with \code{it}, \code{each},
\code{last}, \code{get}) fetches the rows and goes on in memory. After
\code{group-by}, a \code{where} on the aggregates is SQL too (what SQL calls
HAVING), and \kbd{RET} on a group fetches its rows with a query of its own.

\begin{Verbatim}
qss> (sqlite (qss-demo-database) "Orders") (where (> :amount 100))
       (sort-by order-date :desc) (first 20) sql
SELECT *
FROM "Orders"
WHERE ("Amount" > ?)
ORDER BY "OrderDate" IS NULL, "OrderDate" DESC
LIMIT ?
Parameters: (100 20)
\end{Verbatim}

Values never become text in the SQL. Column and table names become keys
mechanically, never by translation: \code{OrderDate} $\to$ \code{order-date},
\code{HT\_Orders} $\to$ \code{ht-orders}, \code{Straße} $\to$ \code{strasse}; the
way back is looked up in the schema. \code{"10kb"} works in SQL too: for
\code{<}, \code{>}, \code{<=}, \code{>=} on every column without a text type,
for \code{=}, \code{/=}, \code{in} only on numeric columns --- so that SQL
returns the same rows as memory.

\section{Storing: \texttt{into}}\label{sec:into}

\code{into} is the way back. The table of a pipe goes into SQLite and comes out
as a query, so the words go on there.

\begin{Verbatim}
qss> (files "~") (where (> :size "10mb")) (into sqlite "big-files")
qss> (sqlite "big-files") (where (> :modified "2026-01-01"))
qss> (nu "sys disks") (into sqlite "disks" :append)
\end{Verbatim}

\begin{itemize}
\item Without \code{:file} the table goes into the central file
  \code{qss-database} (\file{qss.sqlite3} in \code{qss-data-directory},
  by default \file{\textasciitilde/.emacs.d/}); \code{(sqlite "big-files")}
  and \code{(tables)} read from there.
\item Keep the SQLite files on a local disk. If \file{\textasciitilde/.emacs.d}
  lies on a network drive that several Emacs use at once, SQLite's locking
  cannot be trusted there; give each Emacs its own \code{qss-data-directory}.
\item An existing table is not touched unless \code{:replace} or
  \code{:append} says so. All in one transaction.
\item Sizes and durations are stored as INTEGER (bytes, nanoseconds), times as
  ISO text \code{2026-10-22 07:11:05} (\code{qss-time-format}); the types of
  \qss{} are recorded in a table \code{\_qss\_columns}, so sizes show as
  \code{9.8k} again.
\item A file keeps its path, so \kbd{RET} on the stored row opens the file;
  other row types name the fields to keep with \code{qss-stored-fields}.
\end{itemize}

\section{Changing rows}\label{sec:changing}

SQLite is a file to work in, not only to read. A table buffer of an SQLite table
(\code{[buffer]} after a table in the REPL) can be edited:

\noindent\begin{tabularx}{\linewidth}{@{}lX@{}}
\toprule
\kbd{e} & change the cell at point; \kbd{C-u e} makes it empty (NULL) \\
\kbd{d} & mark the row for deletion (again: unmark) \\
\kbd{u} & forget the changes of the row \\
\kbd{C-c C-c} & write all changes, in one transaction \\
\kbd{C-c C-k} & forget all changes \\
\bottomrule
\end{tabularx}

The changes are shown where they are made --- changed cells highlighted, deleted
rows struck through, a count in the mode line --- and written only by
\kbd{C-c C-c}, all or nothing: if one row cannot be written (someone deleted it
meanwhile), none is. \kbd{g} and killing the buffer ask while there are
changes. A numeric column takes numbers, also \code{"10kb"} and \code{"3hr"};
empty input there means NULL.

For many rows at once there are two words; they return how many rows they
changed:

\begin{Verbatim}
qss> (sqlite "big-files") (where (> :size "1gb")) (update-rows :note "archive")
qss> (sqlite (qss-demo-database) "Orders") (where (= :status "cancelled")) (delete-rows)
\end{Verbatim}

Only the rows of the query change: its \code{where}, \code{first} and
\code{sort-by} count. Without \code{where} or \code{first} the words refuse to
touch every row, unless \code{:all-rows t} says so. The values are Lisp and
evaluated.

Rows are found by their rowid. So only queries straight on a table can be
changed --- not views, joins or groups; \qss{} says so instead of guessing.

% ======================================================================
\chapter{Nushell}\label{sec:nushell}

\code{(nu \emph{command})} is a source: Nushell runs the command in
\code{default-directory} and the result becomes a table. In the middle of a pipe
\code{nu} works the other way round as well: the table so far goes to Nushell as
\code{\$in}, and what comes back is a table again (or a single value).

\begin{Verbatim}
qss> (nu "ps | where cpu > 5")
qss> (nu "open export.csv") (where (> :amount 100)) (sort-by date :desc)
qss> (files) (nu "where size > 10kb | get size | math sum")
qss> (sqlite (qss-demo-database) "Orders") (where (is-null :status))
       (nu "to csv | save -f open-orders.csv")
\end{Verbatim}

Both ways use NUON, Nushell's own format, which \qss{} reads and writes itself
(\file{qss-nuon.el}). Unlike JSON it keeps the types: file sizes, durations and
datetimes arrive as columns of the types \code{size}, \code{duration} and
\code{time}, and go back to Nushell as such, so \code{where size > 10kb} works
on a table from Emacs.

% ======================================================================
\chapter{Presentations and accept}\label{chap:present}

\section{The idea}

The idea comes from CLIM, the Common Lisp Interface Manager, and from Dynamic
Windows on the Symbolics Lisp machines before it: \emph{output remembers what it
shows}, so it can serve as input later. In a CLIM listener you type a command,
and where it wants a file you click a file that some earlier output printed;
the program receives the file object, not its name as text.

\qss{} does the same in Emacs buffers. A buffer is the store: whatever is
shown in a visible window can be chosen.

\section{Presentations}

A \emph{presentation} is a piece of text together with an object and a
\emph{presentation type}.

\begin{itemize}
\item Every \textbf{row} of a table is a presentation of its row object: a file
  row of its \code{qss-file} (type \code{file}, or \code{directory}), a process
  row of its process, a row of \code{tables} of a \code{db-table}, a row of data
  of a \code{row} --- or of a \code{file}, if it has a \code{path}.
\item Every \textbf{cell} is a presentation of its value: \code{size},
  \code{duration} and \code{time} from the type of the column, otherwise from
  the value (\code{integer}, \code{number}, \code{string}, \code{time},
  \code{symbol}). Empty cells are none.
\item What the \textbf{REPL prints} --- a number, a line of text, a time --- is
  one too.
\item \textbf{Any text in any buffer}, made with \code{qss-present}.
\end{itemize}

The text stays plain text; the object rides along in the text property
\code{qss-presentation}. Here is a buffer of notes that presents some files; each
name is as good as a file row in a table:

\begin{lstlisting}
(with-current-buffer (get-buffer-create "*to read*")
  (erase-buffer)
  (insert "Files to look at:\n")
  (dolist (f (directory-files "~/notes" t "\\.org\\'"))
    (insert "  ")
    (qss-present (list :path f) 'file (file-name-nondirectory f))
    (insert "\n"))
  (display-buffer (current-buffer)))
\end{lstlisting}

\section{Accept}

\code{(qss-accept \emph{type} \emph{prompt})} asks for an object of a type.

\begin{enumerate}
\item It looks at \textbf{every visible window} --- not at hidden buffers: what
  you see is what you can choose --- and collects the presentations of the type
  or of any type below it.
\item It \textbf{highlights} them; the prompt tells how many: \enquote{with
  file: (3 shown, click one or type)}.
\item You answer by a \textbf{click} on one (or \kbd{RET} on it): the answer is
  the object itself. Or you \textbf{type} its label, with completion over the
  highlighted ones. Or you type something new, and the text is \textbf{parsed}
  for the type: a path for \code{file}, \code{"10kb"} for \code{size},
  \code{"3hr"} for \code{duration}, \code{"2026-10-22"} for \code{time}.
\item \kbd{C-g} gives up. The highlighting is gone either way.
\end{enumerate}

A row whose object fits is offered as a whole row; otherwise its cells are
looked at. So \code{file} offers file rows, \code{value} the cells.

\section{Types}

\begin{Verbatim}
t ─┬─ object ─┬─ file ─ directory
   │          ├─ process, db-table, row
   └─ value ──┬─ number ─ integer ─┬─ size
              │                    └─ duration
              └─ string, time, symbol
\end{Verbatim}

A directory is a file too, a size a number: accepting \code{file} offers
directories as well, accepting \code{directory} only directories. New types come
with \code{qss-define-presentation-type}, with a parser for typed text; a type
without one uses the parser of the nearest type above it. Section~\ref{ref:types}
lists them all.

\section{Examples}

\paragraph{A value into the input.} In the REPL, while typing, \kbd{C-c C-v}
accepts a \code{value}:

\begin{Verbatim}
qss> (sqlite (qss-demo-database) "Orders")
...
    4714           3 2026-10-02   45.0 paid
qss> (where (= :customer-no ▮        C-c C-v: all cells light up; click the 3
qss> (where (= :customer-no 3▮
\end{Verbatim}

A number goes in as a number, text as a string, a time as
\code{"2026-10-22 07:11:05"}. \kbd{C-c C-f} does the same with the name of a
file: \code{!wc -l} \kbd{C-c C-f}, click a file row.

\paragraph{Two files from two places.} Show two directories, each in a buffer of
its own (\code{[buffer]}), side by side. On \file{main.el} in the left window,
\kbd{a} $\to$ \emph{Compare with file\ldots}: the left row is the first file;
every file row in both windows lights up; click \file{main.el} on the right,
and \code{diff} shows the differences.

\paragraph{Only what fits.} On a file, \emph{Copy into directory\ldots} lights up
only directories --- in every visible window, the notes buffer included.

\paragraph{A value as the start.} On a size cell, \kbd{a} offers \emph{Files
larger than this, here}.

\paragraph{From Lisp.}
\begin{lstlisting}
(let ((dir (qss-accept 'directory "Back up into: ")))
  (copy-directory "~/notes" (qss-file-name dir)))
\end{lstlisting}

\section{Commands with typed arguments}\label{sec:commands}

\code{qss-define-command} defines a command whose arguments are accepted by
type, after \code{define-command} in CLIM:

\begin{lstlisting}
(qss-define-command my-open-both "Open side by side"
    ((a file "Open: ") (b file "next to: "))
  "Show the files A and B side by side."
  (find-file (qss-file-name a))
  (find-file-other-window (qss-file-name b)))
\end{lstlisting}

\begin{itemize}
\item \kbd{M-x my-open-both} accepts both files, each by a click or by typing.
\item On every presentation of the type of the first argument --- every file
  row, every presented file in any buffer --- \kbd{a} and the right mouse button
  offer \emph{Open side by side\ldots}. That presentation is the first argument;
  the others are accepted.
\item \code{qss-file-name} turns what was chosen into a file name: a file row, a
  stored row, a typed path, or a row of an own type with a method for it.
\end{itemize}

\section{Compared with CLIM}

\qss{} takes the core --- presentations with types, accept with highlighting,
commands with typed arguments offered on the objects --- and leaves much out:
there are no presentation translators besides the first argument of a command,
no parsing of whole command lines in an input editor, no views, and only visible
windows are searched. What it adds is that every Emacs buffer can take part.

% ======================================================================
\chapter{Contracts}\label{chap:contracts}

\section{Why}

\qss{} is built around generic functions that others --- or you, later ---
fill with methods. Tests check examples; contracts check every call. A contract
says what a function takes and what it returns, and when it is broken, the
message says \emph{who} broke it:

\begin{Verbatim}
qss> (processes) (first -1)
Contract violation: qss-first: argument 2 must satisfy natnump, got -1
  blame: the caller (outside the checked functions)
  contract: (-> qss-tabular (opt natnump) qss-result)
\end{Verbatim}

\begin{Verbatim}
Contract violation: qss-columns: the result must satisfy
(listof qss-column-spec), got (name 42) (element 1 is 42, not qss-column-spec)
  blame: qss-columns itself, the method for my-thing
  contract: (-> any (listof qss-column-spec))
\end{Verbatim}

\section{Blame}

The contracts follow Racket. A contract is an \code{:around} advice on the
function: it sees the arguments before the function runs and the result before
it returns, so it knows the direction of the fault exactly.

\begin{itemize}
\item An argument that breaks the contract is the fault of the
  \textbf{caller}. Every contract binds the name of its function while it runs,
  so if the caller is a checked function itself, it is named:
  \enquote{blame: the caller, qss-continue}. The checked functions are the
  boundaries, as modules are in Racket.
\item A result that breaks it is the fault of the \textbf{function} --- for a
  generic function, of the method for the type of the first argument. This is
  where own row types are caught.
\end{itemize}

\section{The interface}

The whole interface is in one file, \file{qss-interface.el}, readable as the
reference of \qss; section~\ref{ref:contracts} lists it. The protocol of rows
comes first, since methods for own types are where mistakes happen:

\begin{lstlisting}
(qss-define-contract qss-columns (-> any (listof qss-column-spec)))
(qss-define-contract qss-sqlite  (-> stringp (opt (or/c stringp symbolp))
                                     qss-relation-p))
(qss-define-contract qss-accept  (-> qss-presentation-type-name (opt stringp) any))
\end{lstlisting}

\kbd{C-h f} shows the contract of a function, and the contracts are a source:
\code{(contracts) (where (like :function "\%sql\%"))}.

The contracts stand beside the code, not in it. They hold for every method of a
generic function, also methods added later, and survive redefinition. They are
put only on the boundaries: functions that run once per row or per cell
(\code{qss-field}, the formats of values) have none. There they doubled the time
of a pipe over 20000 rows; at the boundaries they cost a few percent.
\code{qss-contract-enabled} turns them all off.

\section{The language}

\noindent\begin{tabularx}{\linewidth}{@{}lX@{}}
\toprule
\code{stringp}, \code{natnump}, \code{qss-table-p} \ldots & a predicate \\
\code{any} & anything \\
\code{(or/c \ldots)}, \code{(and/c \ldots)}, \code{(not/c \emph{c})} & one of, all of, not \\
\code{(one-of \emph{v}\ldots)} & one of these values \\
\code{(listof \emph{c})}, \code{(list/c \ldots)}, \code{(cons/c \emph{a} \emph{d})} & lists and conses \\
\code{(lambda (x) \ldots)} & a predicate written in place \\
\code{(-> \emph{arg}\ldots{} \emph{result})} & a function; \code{(opt \emph{c})} an optional argument (nil is always allowed), \code{(rest \emph{c})} the rest \\
a name & a contract named with \code{qss-define-contract-type} \\
\bottomrule
\end{tabularx}

A function passed as an argument is checked to be a function, but not wrapped.

\section{For your own code}

\file{qss-contract.el} depends on nothing else of \qss{} and can check any Emacs
Lisp:

\begin{lstlisting}
(require 'qss-contract)
(qss-define-contract-type 'my-weight '(and/c numberp (lambda (w) (>= w 0))))
(qss-define-contract my-tare (-> my-weight my-weight my-weight))
\end{lstlisting}

% ======================================================================
\chapter{A row type of your own}\label{chap:own}

This chapter builds a complete row type: the Emacs Lisp libraries on
\code{load-path}. It shows how the parts of \qss{} fit together --- a library
becomes a table row, a presentation of a file, something to store in SQLite, the
argument of a command, and a checked source. The code is the file
\file{doc/qss-manual-example.el}; the tests check that it works.

\section{The row object and the source}

Any Lisp object can be a row. Here a structure, and a function that returns the
table:

\lstinputlisting{example-1.el.txt}
\lstinputlisting{example-2.el.txt}

\section{The protocol}

Three methods make a library a row: its columns, its fields, what \kbd{RET} and
\kbd{a} do. Fields need not be columns: \code{compiled} can be used in
conditions but is not shown.

\lstinputlisting{example-3.el.txt}

Now this works:

\begin{Verbatim}
qss> (libraries "^qss-") (sort-by size :desc) (first 5)
qss> (libraries) (where (> :size "100kb")) count
qss> (libraries "^subr") (where (not (qss-field it 'compiled)))
\end{Verbatim}

The last line uses the field \code{compiled}; because it calls a function, it
is a condition in Lisp with \code{it}, not one that could become SQL.

\section{A library is a file}

A new presentation type below \code{file}, a method that gives it to library
rows, a label, and a method for \code{qss-file-name}:

\lstinputlisting{example-4.el.txt}

With that, everything that takes a \code{file} takes a library: \emph{Compare
with file\ldots} and \emph{Copy into directory\ldots} appear on library rows,
\kbd{C-c C-f} inserts their path, and \code{(qss-accept 'library)} reads a
typed name such as \code{subr}.

\section{Storing and a command of its own}

The path is a column already, so nothing more needs to be stored; a stored row
has a \code{path} and is a file again. A command for libraries only:

\lstinputlisting{example-5.el.txt}
\lstinputlisting{example-6.el.txt}

On a library row, \kbd{a} now offers \emph{Show what it requires}, and
\kbd{M-x my-library-requires} accepts a library by a click or by its name.

\section{A contract}

\lstinputlisting{example-7.el.txt}

\code{(my-libraries 42)} is now reported as the fault of the caller. The
contracts of the protocol check the methods above: had \code{qss-columns}
returned something that is not a list of columns, the message would have named
the method for \code{my-library}.

% ======================================================================
\chapter{A cross reference of Emacs Lisp}\label{chap:xref}

\file{qss-elisp-xref.el} makes a database of all Emacs Lisp it finds: the Lisp
of Emacs itself and the installed packages. Every definition, every use of a
name with what it is there, every call after macro expansion --- for code that
is loaded or not, in files compressed or not. It is the hardest test \qss{} has
had: three million rows, and every question an ordinary pipe.

\section{Building it}

\begin{Verbatim}
M-x qss-elisp-xref-index
\end{Verbatim}

Two steps, each in an Emacs of its own in the background; the messages are in
the buffer \code{*qss-elisp-xref*}. From a shell:

\begin{Verbatim}
emacs --batch -L ~/.emacs.d/lib/qss -l qss-elisp-xref -f qss-elisp-xref-batch [DIR...]
emacs --batch -L ~/.emacs.d/lib/qss -l qss-elisp-xref -f qss-elisp-xref-calls-batch
\end{Verbatim}

Without directories they take the Lisp of the running Emacs and
\code{package-user-dir}. The batch Emacs does not read the init file, so it
does not know a \code{package-user-dir} or \code{qss-data-directory} set
there: give the package directory as an argument, and the file with
\code{-{}-eval '(setq qss-elisp-xref-database "\ldots")'} before \code{-f}, in
both steps. \kbd{M-x qss-elisp-xref-index} passes both on by itself.

\begin{enumerate}
\item \textbf{Reading.} The Lisp reader with the position of every symbol
  (\code{read-positioning-symbols}); nothing is loaded or run. Each top-level
  form is walked: definitions, symbols with their context, strings.
\item \textbf{Compiling.} Each file is compiled into a directory thrown away;
  the byte compiler's call tree gives the calls after macro expansion. A
  \code{when} is gone, the \code{message} inside it stays.
\end{enumerate}

Each file is one transaction, and only one file is in memory at a time; a run
broken off loses nothing. The next run reads only files that changed (by their
SHA-1) and drops files that are gone. A question in a batch run (a coding
system, say) becomes an error of that file instead of waiting for ever.

The database is \file{elisp-xref.sqlite3} in \code{qss-data-directory}.
Over the Lisp of Emacs 29.3 (Debian) it held:

\noindent\begin{tabularx}{\linewidth}{@{}lrX@{}}
\toprule
\textbf{Table} & \textbf{Rows} & \textbf{One row per} \\
\midrule
\code{files} & 1\,626 & file: library, origin (emacs, elpa, user), package and version, Commentary, SHA-1, error \\
\code{forms} & 100\,543 & top-level form, with its source text and its first and last line \\
\code{definitions} & 80\,119 & defun, defmacro, defvar, defcustom, cl-defgeneric, cl-defmethod (with specializers), cl-defstruct, local cl-flet and cl-labels \ldots: arglist, whole docstring, interactive; also the primitives of C \\
\code{occurrences} & 2\,976\,276 & symbol or string (not docstrings, not numbers): line, column, context, the function around it and the argument place \\
\code{calls} & 257\,912 & caller and callee after macro expansion \\
\code{compiled} & 1\,626 & which files were compiled, and why not \\
\bottomrule
\end{tabularx}

Reading took 2\,min 28\,s, compiling 1\,min 45\,s (two cores of a 2.1\,GHz
Xeon, in a virtual machine); the file is 430\,MB. 31 files say
\code{no-byte-compile}; all others compiled.

The context of a symbol tells what it is there: \code{head} (called),
\code{macro-use}, \code{special-form}, \code{function} (\code{\#'f}),
\code{quoted}, \code{defining}, \code{parameter}, \code{binding}, \code{other},
or \code{string}. The strings with their \code{parent\_head} are all the texts a
program shows to people --- a start for a translation.

\section{Asking it from the keyboard}

\code{qss-xref-mode} makes the database a backend of xref, the cross references
built into Emacs. The keys are the usual ones:

\noindent\begin{tabularx}{\linewidth}{@{}lX@{}}
\toprule
\kbd{M-.} & to the definition; a generic function lists its methods with their specializers \\
\kbd{M-?} & every use in code (not in strings, comments or docstrings), each with its line \\
\kbd{C-M-.} & definitions by parts of the name: \code{sample size} finds \code{sample-size} \\
\kbd{M-,} & back \\
\bottomrule
\end{tabularx}

\begin{lstlisting}
(add-hook 'emacs-lisp-mode-hook #'qss-xref-mode)
\end{lstlisting}

Where the database knows nothing --- code not indexed yet, the C source of a
primitive --- Emacs' own Elisp backend answers (\code{qss-xref-fallback}). The
tables of the questions below have the mode on as well, so \kbd{M-.} works on a
name in a row.

The questions are commands, each with the name at point as default (in a qss
table: the value or the row at point), and in one keymap for a prefix of one's
own:

\begin{lstlisting}
(keymap-global-set "C-c x" qss-xref-command-map)
\end{lstlisting}

\noindent\begin{tabularx}{\linewidth}{@{}llX@{}}
\toprule
\textbf{Key} & \textbf{Command} & \textbf{Shows} \\
\midrule
\kbd{u} & \code{qss-xref-uses} & every place a name appears, with the definition it stands in \\
\kbd{c} & \code{qss-xref-callers} & the functions that call it, after macro expansion \\
\kbd{e} & \code{qss-xref-callees} & what a function calls \\
\kbd{d} & \code{qss-xref-definitions} & kind, arglist, line and docstring \\
\kbd{t} & \code{qss-xref-texts} & the texts a library shows to people \\
\kbd{i} & \code{qss-elisp-xref-index} & build or update the database \\
\kbd{m} & \code{qss-xref-mode} & the xref backend on or off \\
\bottomrule
\end{tabularx}

\kbd{RET} on a row opens the file at line and column, compressed or not. The
menu of a row (\kbd{a}, right mouse button) asks the same questions for the
name in it, so one walks from a use to its callers, from there to their
callees, and on.

\section{Asking it in Lisp}

\code{(xref "\emph{table}")} is a source: the table of the database. Every
question is a pipe:

\begin{lstlisting}
;; who calls string-trim, at the line of the caller's definition
(qss-display-> (xref "calls") (where (= :callee "string-trim"))
               (join (xref "definitions") :on ((:file-id :file-id) (:caller :name))
                     :kind left)
               (join (xref "files") :on (:file-id :id))
               (select :library :caller :line :path)
               (sort-by :caller))

;; the macros used most, over all three million rows
(qss-display-> (xref "occurrences") (where (= :context "macro-use"))
               (group-by :value) (sort-by :count :desc) (first 20))

;; who still uses what became obsolete
(qss-display-> (xref "occurrences") (where (in :value '("defadvice" "when-let")))
               (join (xref "files") :on (:file-id :id))
               (group-by :origin :package :value) (sort-by :count :desc))
\end{lstlisting}

The words run as SQL in SQLite; only the rows shown come into Emacs.
\file{qss-elisp-xref-queries.el} has fifteen such questions, one form each, to
try with \kbd{C-x C-e}. The database is an ordinary SQLite file, so
\code{sqlite3} on the command line reads it too.

\section{How fast, and how exact}

Measured over the same Lisp of Emacs 29.3, warm (the files in the page cache),
the best of three runs; the counts are what each finds.

\noindent\begin{tabularx}{\linewidth}{@{}>{\raggedright\arraybackslash}Xrrrr@{}}
\toprule
& \code{seq-reduce} & \code{string-trim} & \code{gamegrid-set-cell} & \code{message} \\
\midrule
\kbd{M-?} on the database & 1\,ms / 26 & 2\,ms / 55 & 1\,ms / 30 & 159\,ms / 5\,747 \\
\code{rg -w}, sources unpacked & 20\,ms / 26 & 21\,ms / 100 & 21\,ms / 30 & 46\,ms / 18\,230 \\
\code{grep -rnw}, sources unpacked & 78\,ms / 26 & 61\,ms / 100 & 67\,ms / 30 & 94\,ms / 18\,230 \\
Emacs' own \kbd{M-?} in the directory, unpacked & 127\,ms / 26 & 86\,ms / 61 & 73\,ms / 30 & 639\,ms / 9\,480 \\
Emacs' own \kbd{M-?}, \file{.el.gz} as installed & 93\,ms / 0 & 12\,ms / 0 & 11\,ms / 0 & 18\,ms / 237 \\
\code{zgrep}, \file{.el.gz} as installed & 3.4\,s / 26 & 3.6\,s / 100 & 3.5\,s / 30 & 3.6\,s / 18\,230 \\
\midrule
\kbd{M-.} on the database & 1\,ms / 1 & 1\,ms / 1 & 1\,ms / 1 & 1\,ms / 1 \\
Emacs' own \kbd{M-.} & 1\,ms / 1 & 0\,ms / 1 & 0\,ms / \textbf{0} & 1\,ms / 1 \\
\midrule
\code{qss-xref-callers} & 8\,ms / 14 & 10\,ms / 42 & 4\,ms / 15 & 401\,ms / 1\,000\footnotemark \\
\bottomrule
\end{tabularx}
\footnotetext{A table shows at most \code{qss-sql-display-limit} rows.}

What the numbers say:

\begin{itemize}
\item \textbf{Fast where it counts.} On unpacked files in the cache, ripgrep
  is fast too. But Emacs installs its Lisp compressed: there Emacs' own
  \kbd{M-?} finds nothing at all, and \code{zgrep} needs three and a half
  seconds per name. The database answers in one or two milliseconds.
\item \textbf{Exact.} For \code{string-trim} grep finds 100 lines: also
  comments, docstrings, and \code{string-trim-left} (a hyphen ends a word for
  grep). The database finds the 55 uses in code, each with what it is there.
  For \code{message}: 5\,747 uses against 18\,230 lines.
\item \textbf{Not only loaded code.} \code{gamegrid-set-cell} is the function
  Tetris draws with. As long as \file{gamegrid} is not loaded, Emacs' own
  \kbd{M-.} does not know it; the database does.
\item \textbf{Questions grep cannot ask.} Callers after macro expansion,
  callees, the methods of a generic function with their specializers, the
  texts of a library, counts over everything: the macros used most, grouped
  over three million rows, in 2.6\,s; the functions with the most texts in
  0.4\,s.
\end{itemize}

% ======================================================================
\chapter{Reference}\label{chap:reference}

This chapter is written from the code by \file{doc/qss-manual-reference.el}.

{\small\sloppy
% Written by qss-manual-reference.el from the code of qss 0.20.0.  Do not edit.
\section{Words}\label{ref:words}
\begin{longtable}{@{}p{0.16\linewidth}p{0.26\linewidth}p{0.52\linewidth}@{}}
\toprule
\textbf{Word} & \textbf{Function} & \textbf{Meaning} \\
\midrule
\endhead
\code{collect} & \code{qss-collect} & Return the rows of TABLE as a Lisp list, at the end of a pipe. \\
\code{columns} & \code{qss-columns-of} & Return the column names of TABLE as a table. \\
\code{count} & \code{qss-count} & Return the number of rows of TABLE. \\
\code{delete-rows} & \code{qss-delete-rows} & Delete the stored rows of TABLE, an SQLite query; return their number. \\
\code{each} & \code{qss-each} & Call FN on every row of TABLE; the results are the new rows. \\
\code{filter} & \code{qss-filter} & Keep the rows of TABLE that satisfy COND (another name of \code{qss-where}). \\
\code{first} & \code{qss-first} & Return the first N rows of TABLE. \\
\code{get} & \code{qss-get} & Return the values of column COL of TABLE as a table with one column. \\
\code{group-by} & \code{qss-group-by} & Group the rows of TABLE by the columns COLS, as Nushell's group-by --to-table. \\
\code{group-sum} & \code{qss-group-sum} & Group the rows of TABLE by the columns in ARGS, with aggregates. \\
\code{head} & \code{qss-head} & Return the first N rows of TABLE (default 10), as \code{qss-first} does. \\
\code{into} & \code{qss-into} & Store TABLE in TARGET (a symbol) and return what was stored, as a source. \\
\code{join} & \code{qss-join} & Join TABLE with OTHER: rows whose join columns are equal, side by side. \\
\code{last} & \code{qss-last} & Return the last N rows of TABLE, in their order. \\
\code{length} & \code{qss-length} & Return the number of rows of TABLE, as \code{qss-count} does. \\
\code{nu} & \code{qss-nu-pipe} & Send TABLE through the Nushell COMMAND and return the result. \\
\code{reject} & \code{qss-reject} & Return TABLE with all columns except COLS (cut down as in \code{qss-select}). \\
\code{select} & \code{qss-select} & Return TABLE with only the columns COLS. \\
\code{sort-by} & \code{qss-sort-by} & Return TABLE sorted by KEYS: columns, each followed by :asc (default) \\
\code{sql} & \code{qss-sql} & Return the query REL as text with its parameters. \\
\code{summarize} & \code{qss-summarize} & Return one row with the count of TABLE and the aggregates AGGS. \\
\code{to} & \code{qss-to} & Return TABLE as text in FORMAT, one of \code{qss-text-formats}. \\
\code{update-rows} & \code{qss-update-rows} & Set columns in the stored rows of TABLE, an SQLite query; return their number. \\
\code{where} & \code{qss-where} & Keep the rows of TABLE that satisfy COND. \\
\bottomrule
\end{longtable}
\section{Sources}\label{ref:sources}
\begin{longtable}{@{}p{0.16\linewidth}p{0.26\linewidth}p{0.52\linewidth}@{}}
\toprule
\textbf{Source} & \textbf{Function} & \textbf{Meaning} \\
\midrule
\endhead
\code{contracts} & \code{qss-contracts} & Return a table of the functions with a contract and their contracts. \\
\code{files} & \code{qss-files} & Return a table of the files in DIR (default \code{default-directory}). \\
\code{lines} & \code{qss-lines} & Return the lines of FILE as a table: columns \code{number} and \code{line}. \\
\code{nu} & \code{qss-nu} & Run the Nushell COMMAND and return its result as a table. \\
\code{odbc} & \code{qss-odbc-source} & Return the table TABLE of the ODBC data source DSN as a query. \\
\code{processes} & \code{qss-processes} & Return a table of the running processes (from \code{process-attributes}). \\
\code{run} & \code{qss-run} & Run PROGRAM with ARGS, without a shell, and return its output as a table. \\
\code{sqlite} & \code{qss-sqlite} & Return the table TABLE of the SQLite file FILE as a query. \\
\code{tables} & \code{qss-tables} & Return the tables of the SQLite file FILE (default \code{qss-database}). \\
\bottomrule
\end{longtable}
\section{Options}\label{ref:options}
All options are in the customization group \code{qss} (\code{M-x customize-group RET qss}).

\begin{longtable}{@{}p{0.34\linewidth}p{0.18\linewidth}p{0.42\linewidth}@{}}
\toprule
\textbf{Option} & \textbf{Default} & \textbf{Meaning} \\
\midrule
\endhead
\code{qss-contract-enabled} & \code{t} & Non-nil means contracts are checked. \\
\code{qss-copy-format} & \code{text} & Format in which \textbackslash{}\code{w} copies a table; see \code{qss-to}. \\
\code{qss-data-directory} & \code{nil} & Directory of the SQLite files of qss; nil means \code{user-emacs-directory}. \\
\code{qss-database} & \code{nil} & The central SQLite file of qss; nil means qss.sqlite3 in \code{qss-data-directory}. \\
\code{qss-max-column-width} & \code{60} & Maximum width of a text column without its own :width, in characters. \\
\code{qss-menu-label-width} & \code{70} & Maximum width of a label in the actions at point. \\
\code{qss-nu-arguments} & \code{("--no-config-file")} & Arguments for \code{qss-nu-program} before -c. \\
\code{qss-nu-format} & \code{nuon} & How data goes between Emacs and Nushell. \\
\code{qss-nu-program} & \code{"nu"} & The Nushell program. \\
\code{qss-nu-size-columns} & \code{(size mem virtual total free used available)} & Columns whose numbers are sizes in bytes, when \code{qss-nu-format} is \code{json}. \\
\code{qss-repl-max-rows} & \code{30} & Maximum number of rows a table shows inside the qss REPL. \\
\code{qss-repl-prompt} & \code{"qss\textgreater{} "} & Prompt of the qss REPL. \\
\code{qss-sql-display-limit} & \code{1000} & Maximum number of rows a query fetches for display. \\
\bottomrule
\end{longtable}
\section{Faces}\label{ref:faces}
\begin{longtable}{@{}p{0.34\linewidth}p{0.62\linewidth}@{}}
\toprule
\textbf{Face} & \textbf{Used for} \\
\midrule
\endhead
\code{qss-accept} & Face of the presentations \code{qss-accept} offers. \\
\code{qss-edit-changed} & Face of a cell changed but not yet written. \\
\code{qss-edit-deleted} & Face of a row marked for deletion. \\
\code{qss-repl-directory} & Face of the directory in the qss REPL prompt. \\
\code{qss-repl-error} & Face of error messages in the qss REPL. \\
\code{qss-repl-prompt} & Face of the qss REPL prompt. \\
\code{qss-repl-title} & Face of the title line above a table in the qss REPL. \\
\code{qss-table} & Face of qss tables. \\
\bottomrule
\end{longtable}
\section{Presentation types}\label{ref:types}
\enquote{Text} says whether typed text can be read as the type.

\begin{longtable}{@{}p{0.15\linewidth}p{0.15\linewidth}p{0.08\linewidth}p{0.54\linewidth}@{}}
\toprule
\textbf{Type} & \textbf{Below} & \textbf{Text} & \textbf{Meaning} \\
\midrule
\endhead
\code{db-table} & \code{object} &  & A table of an SQLite file. \\
\code{directory} & \code{file} & yes & A directory. \\
\code{duration} & \code{integer} & yes & A duration in nanoseconds; \enquote{3hr} can be typed. \\
\code{file} & \code{object} & yes & A file or directory. \\
\code{integer} & \code{number} & yes & A whole number. \\
\code{line} & \code{object} &  & A line of a file (\code{lines}). \\
\code{number} & \code{value} & yes & A number. \\
\code{object} & \code{t} &  & Something with actions: a row object. \\
\code{process} & \code{object} & yes & A system process. \\
\code{row} & \code{object} &  & A row of data (a property list). \\
\code{size} & \code{integer} & yes & A size in bytes; \enquote{10kb} can be typed. \\
\code{string} & \code{value} & yes & Text. \\
\code{symbol} & \code{value} & yes & A symbol. \\
\code{time} & \code{value} & yes & A point in time; \enquote{2026-10-22} or \enquote{2026-10-22 07:30} can be typed. \\
\code{value} & \code{t} & yes & A plain value: number, text, time. \\
\bottomrule
\end{longtable}
\section{Commands with typed arguments}\label{ref:commands}
\begin{longtable}{@{}p{0.27\linewidth}p{0.25\linewidth}p{0.42\linewidth}@{}}
\toprule
\textbf{Command} & \textbf{In the menu} & \textbf{Arguments} \\
\midrule
\endhead
\code{qss-compare-files} & Compare with file & \code{a file}, \code{b file} \\
\code{qss-copy-file-into} & Copy into directory & \code{file file}, \code{dir directory} \\
\code{qss-files-larger-than} & Files larger than this, here & \code{size size} \\
\bottomrule
\end{longtable}
\section{Named contracts}\label{ref:contract-types}
\begin{longtable}{@{}p{0.3\linewidth}p{0.64\linewidth}@{}}
\toprule
\textbf{Named contract} & \textbf{Contract and meaning} \\
\midrule
\endhead
\code{qss-action} & \code{(cons/c stringp functionp)}\newline An action: (LABEL . FUNCTION). \\
\code{qss-aggregate} & \code{(and/c consp (lambda (a) (memq (car a) qss-aggregates)))}\newline An aggregate: (FN COLUMN) or (FN COLUMN :as NAME), FN one of sum avg min max count. \\
\code{qss-column-name} & \code{(or/c symbolp stringp)}\newline A column name: name, :name or \enquote{name}. \\
\code{qss-column-spec} & \code{(or/c symbolp (cons/c symbolp qss-interface--column-plist-p))}\newline A column: a field name, or (NAME . PLIST) with :type, :format, :align, :width. \\
\code{qss-format} & \code{(or/c (one-of text org md csv tsv) stringp)}\newline A text format of \code{qss-to}. \\
\code{qss-presentation-type-name} & \code{qss-presentation-type-p}\newline A presentation type: file, directory, size, value ... or t. \\
\code{qss-result} & \code{(or/c qss-table-p qss-relation-p)}\newline What most words return: a table, or a query on SQLite. \\
\code{qss-tabular} & \code{(or/c qss-table-p qss-relation-p list recordp)}\newline What the words take: a table, a query, a list of rows, or one row (an object). \\
\bottomrule
\end{longtable}
\section{Contracts}\label{ref:contracts}
The 54 functions of the interface with their contracts, as \file{qss-interface.el} defines them.

\begin{longtable}{@{}p{0.3\linewidth}p{0.64\linewidth}@{}}
\toprule
\textbf{Function} & \textbf{Contract} \\
\midrule
\endhead
\code{qss-accept} & \code{(-\textgreater{} qss-presentation-type-name (opt stringp) any)} \\
\code{qss-actions}\newline\emph{generic} & \code{(-\textgreater{} any (listof qss-action))} \\
\code{qss-actions-at-point} & \code{(-\textgreater{} (listof (listof qss-action)))} \\
\code{qss-as-input}\newline\emph{generic} & \code{(-\textgreater{} any qss-result)} \\
\code{qss-as-table}\newline\emph{generic} & \code{(-\textgreater{} any qss-table-p)} \\
\code{qss-collect}\newline\emph{generic} & \code{(-\textgreater{} any any)} \\
\code{qss-columns}\newline\emph{generic} & \code{(-\textgreater{} any (listof qss-column-spec))} \\
\code{qss-columns-of}\newline\emph{generic} & \code{(-\textgreater{} qss-tabular qss-table-p)} \\
\code{qss-commands-for} & \code{(-\textgreater{} any qss-presentation-type-name (listof qss-action))} \\
\code{qss-continue} & \code{(-\textgreater{} qss-tabular consp any)} \\
\code{qss-copy-table-as} & \code{(-\textgreater{} qss-tabular qss-format any)} \\
\code{qss-count}\newline\emph{generic} & \code{(-\textgreater{} qss-tabular natnump)} \\
\code{qss-define-presentation-type} & \code{(-\textgreater{} symbolp qss-presentation-type-name (opt functionp) (opt stringp) symbolp)} \\
\code{qss-delete-rows}\newline\emph{generic} & \code{(-\textgreater{} qss-tabular (rest any) natnump)} \\
\code{qss-demo-database} & \code{(-\textgreater{} stringp)} \\
\code{qss-each}\newline\emph{generic} & \code{(-\textgreater{} qss-tabular functionp qss-result)} \\
\code{qss-file-name}\newline\emph{generic} & \code{(-\textgreater{} any stringp)} \\
\code{qss-files} & \code{(-\textgreater{} (opt stringp) qss-table-p)} \\
\code{qss-first}\newline\emph{generic} & \code{(-\textgreater{} qss-tabular (opt natnump) any)} \\
\code{qss-get}\newline\emph{generic} & \code{(-\textgreater{} qss-tabular qss-column-name any)} \\
\code{qss-group-by}\newline\emph{generic} & \code{(-\textgreater{} qss-tabular (rest qss-column-name) qss-result)} \\
\code{qss-group-sum}\newline\emph{generic} & \code{(-\textgreater{} qss-tabular (rest (or/c qss-column-name qss-aggregate)) qss-result)} \\
\code{qss-into}\newline\emph{generic} & \code{(-\textgreater{} qss-tabular symbolp (rest any) any)} \\
\code{qss-join}\newline\emph{generic} & \code{(-\textgreater{} qss-tabular qss-tabular (rest any) qss-result)} \\
\code{qss-last}\newline\emph{generic} & \code{(-\textgreater{} qss-tabular (opt natnump) any)} \\
\code{qss-make-table} & \code{(-\textgreater{} stringp list (opt (listof qss-column-spec)) (opt functionp) qss-table-p)} \\
\code{qss-nu} & \code{(-\textgreater{} stringp any)} \\
\code{qss-nu-pipe} & \code{(-\textgreater{} qss-tabular stringp any)} \\
\code{qss-nuon-read} & \code{(-\textgreater{} stringp any)} \\
\code{qss-odbc-source} & \code{(-\textgreater{} stringp (opt (or/c stringp symbolp)) (rest any) (or/c qss-relation-p qss-table-p))} \\
\code{qss-present} & \code{(-\textgreater{} any (opt qss-presentation-type-name) (opt stringp) any)} \\
\code{qss-presentation-label}\newline\emph{generic} & \code{(-\textgreater{} any stringp)} \\
\code{qss-presentation-parse} & \code{(-\textgreater{} qss-presentation-type-name stringp any)} \\
\code{qss-presentation-subtype-p} & \code{(-\textgreater{} qss-presentation-type-name qss-presentation-type-name any)} \\
\code{qss-presentation-type}\newline\emph{generic} & \code{(-\textgreater{} any qss-presentation-type-name)} \\
\code{qss-processes} & \code{(-\textgreater{} qss-table-p)} \\
\code{qss-reject}\newline\emph{generic} & \code{(-\textgreater{} qss-tabular (rest qss-column-name) qss-result)} \\
\code{qss-run} & \code{(-\textgreater{} stringp (rest stringp) qss-table-p)} \\
\code{qss-select}\newline\emph{generic} & \code{(-\textgreater{} qss-tabular (rest qss-column-name) qss-result)} \\
\code{qss-show} & \code{(-\textgreater{} qss-table-p (opt any) bufferp)} \\
\code{qss-sort-by}\newline\emph{generic} & \code{(-\textgreater{} qss-tabular (rest qss-column-name) qss-result)} \\
\code{qss-sql} & \code{(-\textgreater{} any stringp)} \\
\code{qss-sql-name-key} & \code{(-\textgreater{} stringp symbolp)} \\
\code{qss-sql-relation} & \code{(-\textgreater{} any (or/c stringp symbolp) qss-relation-p)} \\
\code{qss-sql-rowid} & \code{(-\textgreater{} any (or/c null integerp list))} \\
\code{qss-sql-schema} & \code{(-\textgreater{} stringp (listof (cons/c stringp list)))} \\
\code{qss-sql-write} & \code{(-\textgreater{} qss-table-p stringp stringp (opt (one-of replace append)) natnump)} \\
\code{qss-sqlite} & \code{(-\textgreater{} stringp (opt (or/c stringp symbolp)) qss-relation-p)} \\
\code{qss-stored-fields}\newline\emph{generic} & \code{(-\textgreater{} any (listof symbolp))} \\
\code{qss-summarize}\newline\emph{generic} & \code{(-\textgreater{} qss-tabular (rest qss-aggregate) qss-result)} \\
\code{qss-tables} & \code{(-\textgreater{} (opt stringp) qss-table-p)} \\
\code{qss-to}\newline\emph{generic} & \code{(-\textgreater{} qss-tabular qss-format stringp)} \\
\code{qss-update-rows}\newline\emph{generic} & \code{(-\textgreater{} qss-tabular (rest any) natnump)} \\
\code{qss-where}\newline\emph{generic} & \code{(-\textgreater{} qss-tabular (or/c qss-condition-p functionp) qss-result)} \\
\bottomrule
\end{longtable}

}

% ======================================================================
\chapter{Where the ideas come from}

Little in \qss{} is new; it puts old ideas together in Emacs.

\begin{description}
\item[Nushell] Data as tables, not text; one vocabulary for every source; the
  pipe. The names of the words, the units \code{10kb} and \code{3hr}, and NUON
  come from there. \url{https://www.nushell.sh/}
\item[dbplyr] A pipe over a database table stays a query and becomes SQL only
  when the rows are needed; the same words work in memory and in the database.
  \url{https://dbplyr.tidyverse.org/}
\item[SQL] The meaning of the comparisons: \code{like}, \code{in}, and NULL as
  unknown in a three-valued logic.
\item[SQLite] One file as the store of everything, source and target alike.
\item[CLIM and Dynamic Windows] Presentations, presentation types,
  \code{accept}, commands with typed arguments offered on the objects shown.
  McCLIM is a free implementation: \url{https://mcclim.common-lisp.dev/}.
\item[CLOS] Generic functions with methods per type (\code{cl-generic} in Emacs):
  the protocol of rows, so a file stays a file and own types fit in with a few
  methods.
\item[Racket contracts] Contracts at the boundaries, and blame: the caller or
  the function. R.\,B.~Findler, M.~Felleisen, \emph{Contracts for Higher-Order
  Functions}, ICFP 2002.
  \url{https://docs.racket-lang.org/guide/contracts.html}
\item[Emacs] \code{vtable} for the tables, the built-in \code{sqlite}, advice
  for the contracts, text properties for the presentations, IELM for the shape
  of the REPL --- and the conviction that the editor is the place to work in.
\end{description}

\appendix

\chapter{Files and tests}

\noindent\begin{tabularx}{\linewidth}{@{}lX@{}}
\toprule
\file{qss.el} & package headers, loads everything, the commands to run the tests \\
\file{qss-table.el} & tables, display, the protocol of rows, \code{files}, \code{processes}, \code{run} \\
\file{qss-pipe.el} & the words, conditions, \code{qss->}, \code{to}, \code{into}, groups, joins \\
\file{qss-sql.el} & SQLite as a source and a store; joins and groups as SQL \\
\file{qss-nu.el}, \file{qss-nuon.el} & Nushell; NUON reader and writer \\
\file{qss-present.el} & presentation types, \code{qss-present}, \code{qss-accept}, commands \\
\file{qss-menu.el} & actions at point, \kbd{a}, the right mouse button, \kbd{w} \\
\file{qss-edit.el} & changing rows of SQLite tables; \code{update-rows}, \code{delete-rows} \\
\file{qss-repl.el} & \kbd{M-x qss} \\
\file{qss-contract.el} & contracts with blame (stands alone) \\
\file{qss-interface.el} & the interface of \qss{} with its contracts \\
\file{qss-examples.el} & \kbd{M-x qss-examples} \\
\file{doc/} & this manual, its example and the writer of its reference \\
\bottomrule
\end{tabularx}

\kbd{M-x qss-run-tests} runs the tests in Emacs; from a shell:

\begin{Verbatim}
emacs --batch -L /path/to/qss -l qss -f qss-run-tests-batch
\end{Verbatim}

The tests run with the contracts on, so every test checks the code against the
contracts as well. The SQLite tests are skipped without SQLite, the Nushell
tests without \code{nu}.

\end{document}
